\title{
MADRS-S
}
\section*{(självskattningsskala)}
\section*{Namn}
\section*{Datum}
Genom att besvara följande nio frågor kan du och din läkare få en detaljerad bild av hur du mår och om du har symtom, som är typiska för depression. Genom att lägga ihop den "poäng" du får på frågorna får du och din läkare en bild av graden av depression. Sätt en ring runt siffran som du tycker bäst stämmer med hur du mått de senaste tre dagarna. Använd gärma mellanliggande alternativ. Tänk inte alltför länge, utan forsök arbeta snabbt.
\section*{1. Sinnestämning}
Här ber vi dig beskriva din sinnesstämning, om du känner dig ledsen, tungsint eller dyster till mods. Tänk efter hur du har känt dig de senaste tre dagarna, om du har skiftat i humöret eller om det varit i stort sett detsamma hela tiden, och försök särskilt komma ihg om du har känt dig lättare till sinnes om det har känt något positivt.
0 Jag kan känna mig glad eller ledens, alltefter omständigheterna.
1
2 Jag känner mig nedstämd för det mesta, men ibland kan det kännas lättare.
4 Jag känner mig genomgående nedestämd och dyster. Jag kan inte glädja mig åt sådant som vanligen skulle göra mig glad.
5
6 Jag är så totalt nedstämd och olycklig att jag inte kan tänka mig värre.
\section*{2. Oroskänslor}
Här ber vi dig markera i vilken utsträckning du haft känslor av inre spänning, lout och ångest eller odefinierad rädsla under de senaste tre dagarna. Tänk särskilt på hur intensiva känslorna varit, och om de kommit och gått eller funnits nästan hela tiden.
0 Jag känner mig mestadels lung.
2 libland har jag obehagliga känslor av inre oro.
3
4 Jag har ofta en känsla av inre oro som ibland kan bli mycket stark, och som jag måste anstränga mig för att bemästra.
5
6 Jag har fruktansvärda, långvariga eller outhärdliga ångestkänslor.
\section*{3. Sömn}
Här ber vi dig beskriva hur bra du sover. Tänk efter hur länge du sovit och hur god sömnen varit under de senaste tre nätterna. Bedömningen skall avse hur du faktiskt sovit, oavsett om du tagit sömnmedel eller ej. Om du sover mer än vanligt, sätt din markering vid 0
0 Jag sover lugnt och bra och tillräckligt länge för mina behov. Jag har inga särskilda svårigheter att somna.
2 Jag har vissa sömnsvårigheter. Ibland har jag svårt att somna eller sover yltigare eller olrigare än vanligt.
4 Jag sover minst två timarm mindre per natt än normalt. Jag vaknar ofta under natten, även om jag inte blir stård.
6 Jag sover mycket dåligt, inte mer än 2-3 timmar per natt.
\section*{4. Matlust}
Här ber vi dig ta ställning till hur din aptit är, och tänka efter om den på något sätt skilt sig från vad som år normalt för dig. Om du skulle ha bättre aptit än normalt, sätt din markering vid 0
0 Min aptit är som den brukar vara.
2 Min aptit är sämre än vanligt.
3 Min aptit har nästan helt försvunnit. Maten smakar inte och jag måste tvinga mig att åta.
6 Jag vill inte ha någon mat. Om jag skulle få någonting i mig, måste jag övertalas att åta.

\section*{5. Koncentrationsförmåga}
Här ber vi dig ta ställning till din förmåga att hålla tankarna samlade och koncentrera dig på olika aktiviteter. Tänk igenom hur du fungerar vid olika sysslor som kråver olika grad av koncentrationsförmåga, te \(1 \mathrm{ælsning} \mathrm{av}\) komplicerad text, lätt tidningstext och TV-tittande.
0 Jag har inga koncentrationsvärgheter
1 maxentingertilst värt att hålla tankarna samlade på
sådant som normalt skulle fågna im uppmärksamhet (t ex länsing eller TV-tittande).
4 Jag har påtagligt svårt att koncentrera mig på sådant som normalt inte kräver någon ansträngning från min sida (t ex länsing eller samtal med andra människor).
6 Jag kan överhuvudtaget inte koncentrera mig på någonting.
\section*{6. Initiativförmåga}
Här ber vi dig förösöka värdera din handlingskraft. Frågan gäller om du har lätt eller svårt för dig att komma igång med sådant du tycker du bör göra, och i vilken utsträckning du måste övervinna ett inre motstånd när skall ta itu med något.
0 Jag har inga svårigheter med att ta itu med nya uppgifter.
2 När skall jag ta itu med något, tar det emot på ett sätt som inte är normalt för mig.
3 Det krävs en stor ansträngning för mig att ens komma igång med enkla uppgifter som jag vanligtvis utför mer eller mindre rutinmässigt.
6 Jag kan inte förmå mig att ta itu med de enklaste vardagssysslor.
\section*{7. Känslomässigt engagemang}
Här ber vi dig ta ställning till hur du upplever ditt intresse för omvärlden och för andra människor, och för sådana aktiviteter som brukar bereda dig nöje och glädje.
0 Jag är intresserad av omvärlden och engagerar mig i den, och det bereder mig både nöje och glädje.
2 Jag känner mindre starkt för sådant som brukar engagera mig. Jag har svårare än vanligt att bli glad eller svårare att bli arg när det är befogat.
3 Jag kan inte känna något intresse för omvärlden, inte ens för vänner och bekanta.
5
6 Jag har slutat uppleva några känslor. Jag känner mig smärtsamt likgiltig även för mina närmaste.
\section*{8. Pessimism}
Frågan gäller hur du ser på din egen framtid och hur du uppfattard itt eget värde. Tänk efter i vilken utsträckning duer gär självförebräelser, om du plågas av skuldänslor, och om du oroat dig oftere än vanligt för te \(x\) end ienoki eller din hälsa.
0 Jag ser på framtiden med tillförsikt. Jag är på det hela taget ganska nöjd med mig själv.
1 bland klandrar jag mig själv och tycker jag är mindre värd än andra.
4 Jag grubblar ofta över mina misslyckanden och känner mig mindervärdig eller dålig, även om andra tycker annorlunda.
6 Jag ser allting i svart och kan inte se någon ljusning. Det känns som om jag var en alltigenom dålig månniska, och som om jag radlig skulle kunna få någon förlåtesle för det hemska jag gjort.
\section*{9. Livslust}
Frågan gäller din livslust, och om du känt livsleda. Har du tankar på självmord, och i så fall, i vilken utsträckning upplever du detta som en verkligt utväg?
0 Jag har normal apitt på livet.
2 Livet känns inte särskilt meningsfullt men jag önskar ändå inte att jag vore død.
3 Jag tycker ofta det vore bättre att vara död, och trots att jag egentligen inte önskar det, kan självmord ibland kännas som en möjlig väg.
6 Jag är egentligen övertygad om att min enda utväg är att dö, och jag tänker mycket på hur jag bäst skall gå tillväga för att ta mitt eget liv.

Lägg samman poängen från båda sidor av formulåret och ange summan i rutan